\chapter{Vector autoregressions (VAR)}
\label{cap:var}

\section{The model}
A VAR($p$) for a vector of $k$ variables is
\[
y_t=c+A_1y_{t-1}+\dots+A_py_{t-p}+u_t,\qquad \E(u_t)=0,\ \ \E(u_tu_t')=\Sigma.
\]
Each variable is explained by the lags of \emph{all} variables, including its own. Estimation by OLS, equation by equation, is consistent and, when all equations have the same regressors, is equivalent to maximum likelihood estimation (under normality).

\textbf{Stability (stationarity)}: all eigenvalues of the companion matrix must have modulus less than 1. This is the multivariate analogue of the condition $|\phi|<1$.

\section{Choosing the number of lags}
Estimate VARs with $p=1,2,\dots,p_{\max}$ and compare AIC, BIC (or Hannan--Quinn). Check that the residuals have no autocorrelation.

\section{Granger causality}
$x$ \emph{Granger-causes} $y$ if the lags of $x$ help predict $y$ beyond the lags of $y$ itself. It is tested with a joint F test of the coefficients of $x$ in the equation for $y$.

\begin{atencao}[Pitfall: Granger is not causality]
It is a notion of \emph{predictive precedence}. If an omitted third factor moves both variables with different lags, ``Granger causality'' appears with no cause at all. The name is misleading.
\end{atencao}

\begin{simulacao}[Checked by simulation]
I simulated a bivariate VAR(1) with $A=\begin{pmatrix}0.5&0.2\\0&0.4\end{pmatrix}$ ($T=1500$). Estimate: $\hat A=\begin{pmatrix}0.51&0.195\\-0.005&0.399\end{pmatrix}$. The eigenvalues of $A$ have moduli $0.5$ and $0.4$ (stable). The Granger test correctly identified: $y_2$ helps predict $y_1$ (p-value $\approx10^{-15}$) and $y_1$ does not help predict $y_2$ (p-value $0.82$).
\end{simulacao}

\section{Impulse response functions and variance decomposition}
The \textbf{impulse response} shows how each variable reacts, over time, to a shock in one equation. The problem: the errors $u_t$ are usually correlated across equations, so ``a shock only in $y_1$'' does not exist without an additional assumption. The most common solution is the Cholesky decomposition, which imposes a recursive order (the first variable affects the others in the same period, but not the reverse).

\begin{atencao}
\textbf{The order of the variables in the Cholesky decomposition changes the results} when the errors are correlated. This is an identification choice, and it must be justified economically, not decided for convenience. The forecast error variance decomposition (FEVD) inherits the same problem.
\end{atencao}

\begin{exercicio}
A VAR(1) has $A=\begin{pmatrix}0.6&0.3\\0.2&0.5\end{pmatrix}$. Is it stable?
\end{exercicio}
\sol{Eigenvalues: $\lambda^2-(0.6+0.5)\lambda+(0.6\cdot0.5-0.3\cdot0.2)=\lambda^2-1.1\lambda+0.24=0$. So $\lambda=(1.1\pm\sqrt{1.21-0.96})/2=(1.1\pm0.5)/2$, that is, $0.8$ and $0.3$. Both have modulus less than 1: the VAR is stable.}
